\begin{table}[t]
    \centering
    \caption{Hyperparameters introduced in \wmrlshort{}. \emph{Value} represents value used in ALFWorld, $\tau^2$ Bench, respectively.}
    \label{tab:rwml_hyperparams}
    \begin{tabular}{l l l p{4.2cm} p{5.2cm}}
        \toprule
        Stage & Name & Value & Heuristics & Intuition \\
        \midrule
        \wmrlshort{} Data
        & $\tau_d$
        & 0.1, 0.15
        & Set such that ``too easy'' samples correspond to $\sim$30\% of the original dataset
        & Spending more training on medium-to-hard samples better incentivizes the model to learn non-trivial world-model knowledge. \\
        
        \wmrlshort{} Data
        & $\tau_{\mathrm{easy}}$
        & 0.0, 0.0
        & Fixed
        & Spending more training on medium-to-hard samples better incentivizes the model to learn non-trivial world-model knowledge. \\
        
        \wmrlshort{} Data
        & $p$
        & 0.1, 0.1
        & Fixed
        & The final dataset retains mostly medium-to-hard examples while preserving sufficient dataset diversity. \\
        \midrule
        
        \wmrlshort{} Training
        & $\tau_d$
        & 0.2, 0.4
        & Add 0.15 from $\tau_d$ used in data collection, then rounded to nearest 0.2
        & Slightly higher than that of data construction since it uses a fine-tuned model. Values that are too high (\eg{} >0.5) make the world-modeling task overly easy, encouraging generic next-state descriptions. Values that are too low (\eg{} <0.1) approximate exact matching and become overly strict in many settings. \\
        \bottomrule
    \end{tabular}
\end{table}